\let\negmedspace\undefined
\let\negthickspace\undefined
\documentclass[journal]{IEEEtran}
\usepackage[a5paper, margin=10mm, onecolumn]{geometry}
%\usepackage{lmodern} % Ensure lmodern is loaded for pdflatex
\usepackage{tfrupee} % Include tfrupee package

\setlength{\headheight}{1cm} % Set the height of the header box
\setlength{\headsep}{0mm}     % Set the distance between the header box and the top of the text

\usepackage{gvv-book}
\usepackage{gvv}
\usepackage{cite}
\usepackage{amsmath,amssymb,amsfonts,amsthm}
\usepackage{algorithmic}
\usepackage{graphicx}
\usepackage{textcomp}
\usepackage{xcolor}
\usepackage{txfonts}
\usepackage{listings}
\usepackage{enumitem}
\usepackage{mathtools}
\usepackage{gensymb}
\usepackage{comment}
\usepackage[breaklinks=true]{hyperref}
\usepackage{tkz-euclide} 
\usepackage{listings}
% \usepackage{gvv}                                        
\def\inputGnumericTable{}                                 
\usepackage[latin1]{inputenc}                                
\usepackage{color}                                            
\usepackage{array}                                            
\usepackage{longtable}                                       
\usepackage{calc}                                             
\usepackage{multirow}                                         
\usepackage{hhline}                                           
\usepackage{ifthen}                                           
\usepackage{lscape}
\begin{document}

\bibliographystyle{IEEEtran}
\vspace{3cm}

\title{5.2.14}
\author{EE25BTECH11011 - Navya}
% \maketitle
% \newpage
% \bigskip
{\let\newpage\relax\maketitle}

\renewcommand{\thefigure}{\theenumi}
\renewcommand{\thetable}{\theenumi}
\setlength{\intextsep}{10pt} % Space between text and floats
\textbf{Question}:\\
 Solve the following system of linear equations.
 \begin{align}
 x-3y-3=0\\
 3x-9y-2=0
 \end{align}
\textbf{Solution}:\\
The given equations of lines are,
\begin{align}
 \myvec{1 & -3}\vec{x} = 3 \\ 
 \myvec{3 & -9}\vec{x} = 2
\end{align}

In matrix form,
\begin{align}
\myvec{1 & -3 \\ 3 & -9}\vec{x} = \myvec{3 \\ 2}
\end{align}

So the augmented matrix is
\begin{align}
\augvec{2}{1}{1 & -3 & 3 \\ 3 & -9 & 2}
\end{align}

Performing row operations:
\begin{align}
R_2 \to R_2 - 3R_1 
\implies \augvec{2}{1}{1 & -3 & 3 \\ 0 & 0 & -7}
\end{align}

This corresponds to the equation
\begin{align}
0x + 0y = -7,
\end{align}
which is a contradiction.\\

Hence, the system has \textbf{no solution}. The two equations can be written as
\begin{align}
x - 3y = 3, \quad x - 3y = \tfrac{2}{3},
\end{align}
which represent two parallel lines with the same slope but different intercepts. \\


Therefore,The given system of equations is \textbf{inconsistent}, and the two lines are parallel.
\begin{figure}[h!]
    \centering
    \includegraphics[width=0.8\textwidth]{figs/fig7.png}
    \caption{Parallel lines}
    \label{fig:division}
\end{figure}
 
\end{document}